\section{4.11.Evaluación experimental}\label{evaluaciuxf3n-experimental}

\subsection{4.11. Flujo completo del
análisis}\label{flujo-completo-del-anuxe1lisis}

Una vez descritos los distintos módulos que componen EVMAudit, resulta
posible analizar el funcionamiento global de la solución.

El proceso completo implementado por la herramienta se compone de varias
fases consecutivas que permiten transformar un contrato inteligente en
un informe consolidado de vulnerabilidades.

\subsubsection{4.11.1. Recepción del
contrato}\label{recepciuxf3n-del-contrato}

El proceso comienza con la recepción del contrato inteligente que se
desea analizar.

El contrato puede proceder de distintos orígenes:

\begin{itemize}
\tightlist
\item
  ejecución directa de la librería
\item
  aplicación web desarrollada
\item
  futuras integraciones externas.
\end{itemize}

Tras recibir el contrato, EVMAudit detecta automáticamente el nombre del
contrato y configura la versión adecuada del compilador Solidity.

\subsubsection{4.11.2. Obtención de resultados mediante Slither y
Mythril}\label{obtenciuxf3n-de-resultados-mediante-slither-y-mythril}

La siguiente fase consiste en ejecutar las herramientas de análisis
principales.

En primer lugar se ejecuta Slither, aprovechando las ventajas del
análisis estático.

Posteriormente se ejecuta Mythril, incorporando las capacidades
derivadas de la ejecución simbólica.

La combinación de ambas técnicas permite obtener una cobertura superior
a la proporcionada por cada herramienta de forma individual.

!TODO: referencia cruzada a la Sección 2.X (herramientas de análisis).

\subsubsection{4.11.3. Normalización de
resultados}\label{normalizaciuxf3n-de-resultados}

Las salidas generadas por las herramientas se transforman posteriormente
a una estructura común.

Este proceso elimina las diferencias existentes entre ambas herramientas
y permite trabajar con una representación homogénea de las
vulnerabilidades.

La normalización constituye el paso previo necesario para aplicar el
mecanismo de correlación.

\subsubsection{4.11.4. Correlación de
vulnerabilidades}\label{correlaciuxf3n-de-vulnerabilidades}

Una vez normalizados los resultados, EVMAudit agrupa aquellas
vulnerabilidades equivalentes detectadas por varias herramientas.

La correlación permite:

\begin{itemize}
\tightlist
\item
  reducir duplicidades
\item
  aumentar la confianza de determinados hallazgos
\item
  simplificar la interpretación del informe final.
\end{itemize}

Esta fase representa uno de los principales elementos diferenciadores de
la solución desarrollada.

\subsubsection{4.11.5. Generación de propiedades y
fuzzing}\label{generaciuxf3n-de-propiedades-y-fuzzing}

A partir de las vulnerabilidades correlacionadas, el sistema consulta el
catálogo SWC y genera automáticamente las propiedades necesarias para
ejecutar Echidna.

Posteriormente se construye un wrapper Solidity específico y se realiza
una fase adicional de validación mediante fuzzing.

Este proceso proporciona información complementaria sobre determinadas
vulnerabilidades detectadas durante las fases anteriores.

\subsubsection{4.11.6. Generación del informe
final}\label{generaciuxf3n-del-informe-final}

Finalmente, toda la información obtenida durante el análisis es
consolidada en un informe único.

El informe incluye:

\begin{itemize}
\tightlist
\item
  vulnerabilidades detectadas
\item
  herramientas que las han identificado
\item
  nivel de confianza asociado
\item
  resultados del fuzzing
\item
  metadatos adicionales.
\end{itemize}

De este modo, el auditor dispone de una visión unificada del estado de
seguridad del contrato analizado.

\subsubsection{4.11.7. Resumen del proceso}\label{resumen-del-proceso}

El flujo general implementado por EVMAudit puede resumirse en las
siguientes etapas:

\begin{enumerate}
\def\labelenumi{\arabic{enumi}.}
\tightlist
\item
  Recepción del contrato.
\item
  Configuración del entorno.
\item
  Ejecución de Slither.
\item
  Ejecución de Mythril.
\item
  Normalización.
\item
  Correlación.
\item
  Consulta del catálogo SWC.
\item
  Generación del wrapper.
\item
  Ejecución de Echidna.
\item
  Generación del informe.
\end{enumerate}

La secuencia completa permite combinar distintas técnicas de análisis
dentro de una arquitectura unificada y automatizada.

!TODO: insertar diagrama de secuencia del pipeline completo.

!TODO: insertar diagrama de flujo del proceso.

!TODO: insertar referencia a la Figura X.

!TODO: insertar referencia cruzada a las secciones 4.3 a 4.9.

\subsubsection{4.11.X. Evaluación del contrato con una vulnerabilidad
conocida: {[}NOMBRE DEL
CONTRATO{]}}\label{x.-evaluaciuxf3n-del-contrato-con-una-vulnerabilidad-conocida-nombre-del-contrato}

\paragraph{Descripción del contrato}\label{descripciuxf3n-del-contrato}

El contrato \textbf{{[}NOMBRE DEL CONTRATO{]}} fue seleccionado debido a
que presenta una vulnerabilidad conocida relacionada con \textbf{{[}TIPO
DE VULNERABILIDAD{]}}.

Este contrato pertenece a \textbf{{[}FUENTE DEL CONTRATO{]}} y fue
incluido en la evaluación con el objetivo de analizar el comportamiento
de EVMAudit frente a escenarios representativos.

La vulnerabilidad esperada en este caso corresponde a:

\begin{itemize}
\tightlist
\item
  TODO.
\end{itemize}

!TODO: referencia cruzada a la Sección 2.X (descripción de la
vulnerabilidad).

\begin{center}\rule{0.5\linewidth}{0.5pt}\end{center}

\paragraph{Resultados obtenidos por las
herramientas}\label{resultados-obtenidos-por-las-herramientas}

En primer lugar, el contrato fue procesado individualmente por Slither y
Mythril.

Los resultados obtenidos se resumen en la Tabla X.

{\def\LTcaptype{none} % do not increment counter
\begin{longtable}[]{@{}ll@{}}
\toprule\noalign{}
Herramienta & Vulnerabilidades detectadas \\
\midrule\noalign{}
\endhead
\bottomrule\noalign{}
\endlastfoot
Slither & TODO \\
Mythril & TODO \\
\end{longtable}
}

!TODO: insertar referencia a la Tabla X.

\begin{center}\rule{0.5\linewidth}{0.5pt}\end{center}

\paragraph{Resultados tras la
correlación}\label{resultados-tras-la-correlaciuxf3n}

Una vez aplicados los mecanismos de normalización y correlación
implementados por EVMAudit, se obtuvo un total de \textbf{TODO}
vulnerabilidades consolidadas.

La Tabla X muestra la relación entre los hallazgos originales y los
resultados finales.

{\def\LTcaptype{none} % do not increment counter
\begin{longtable}[]{@{}ll@{}}
\toprule\noalign{}
Métrica & Valor \\
\midrule\noalign{}
\endhead
\bottomrule\noalign{}
\endlastfoot
Hallazgos originales & TODO \\
Hallazgos correlacionados & TODO \\
Vulnerabilidades confirmadas & TODO \\
Reducción de duplicidades & TODO \\
\end{longtable}
}

!TODO: insertar referencia a la Tabla X.

\begin{center}\rule{0.5\linewidth}{0.5pt}\end{center}

\paragraph{Resultados del fuzzing mediante
Echidna}\label{resultados-del-fuzzing-mediante-echidna}

A partir de las vulnerabilidades obtenidas se generaron automáticamente
las propiedades correspondientes para Echidna.

Los resultados obtenidos fueron los siguientes:

\begin{itemize}
\tightlist
\item
  Propiedades generadas: TODO.
\item
  Propiedades ejecutadas correctamente: TODO.
\item
  Vulnerabilidades verificadas: TODO.
\item
  Limitaciones observadas: TODO.
\end{itemize}

!TODO: insertar captura o fragmento representativo del resultado de
Echidna.

\begin{center}\rule{0.5\linewidth}{0.5pt}\end{center}

\paragraph{Discusión}\label{discusiuxf3n}

Los resultados obtenidos muestran que \textbf{{[}OBSERVACIÓN
PRINCIPAL{]}}.

En particular:

\begin{itemize}
\tightlist
\item
  TODO.
\item
  TODO.
\item
  TODO.
\end{itemize}

Este caso pone de manifiesto \textbf{{[}CONCLUSIÓN PRINCIPAL DEL
CONTRATO{]}}.

\subsubsection{4.11.X. Evaluación del contrato con varias
vulnerabilidades: {[}NOMBRE DEL
CONTRATO{]}}\label{x.-evaluaciuxf3n-del-contrato-con-varias-vulnerabilidades-nombre-del-contrato}

\paragraph{Descripción del
contrato}\label{descripciuxf3n-del-contrato-1}

El contrato \textbf{{[}NOMBRE{]}} incorpora una vulnerabilidad
relacionada con \textbf{{[}TIPO{]}}, lo que permite estudiar la
capacidad de detección de las herramientas integradas cuando existen
varias debilidades simultáneas.

!TODO: describir brevemente el funcionamiento del contrato.

\begin{center}\rule{0.5\linewidth}{0.5pt}\end{center}

\paragraph{Comparativa entre
herramientas}\label{comparativa-entre-herramientas}

La Tabla X resume las vulnerabilidades detectadas por cada herramienta.

{\def\LTcaptype{none} % do not increment counter
\begin{longtable}[]{@{}lcc@{}}
\toprule\noalign{}
Vulnerabilidad & Slither & Mythril \\
\midrule\noalign{}
\endhead
\bottomrule\noalign{}
\endlastfoot
TODO & {[}X{]} & {[}X{]} \\
TODO & {[}X{]} & {[}-{]} \\
TODO & {[}-{]} & {[}X{]} \\
\end{longtable}
}

Los resultados muestran que ambas herramientas presentan comportamientos
complementarios.

\begin{center}\rule{0.5\linewidth}{0.5pt}\end{center}

\paragraph{Efecto del mecanismo de
correlación}\label{efecto-del-mecanismo-de-correlaciuxf3n}

La aplicación del proceso de correlación permitió:

\begin{itemize}
\tightlist
\item
  eliminar duplicidades;
\item
  unificar nomenclaturas;
\item
  aumentar la confianza de determinados hallazgos.
\end{itemize}

!TODO: insertar número real de vulnerabilidades correlacionadas.

\begin{center}\rule{0.5\linewidth}{0.5pt}\end{center}

\paragraph{Resultados del fuzzing}\label{resultados-del-fuzzing}

El proceso de generación automática produjo un total de \textbf{TODO}
propiedades.

Durante la ejecución de Echidna se observó:

\begin{itemize}
\tightlist
\item
  TODO.
\end{itemize}

!TODO: describir vulnerabilidades no testables.

\begin{center}\rule{0.5\linewidth}{0.5pt}\end{center}

\paragraph{Discusión}\label{discusiuxf3n-1}

Este contrato evidencia que las distintas técnicas de análisis empleadas
por EVMAudit permiten obtener información complementaria y mejorar la
cobertura respecto al uso individual de cada herramienta.

\subsubsection{4.11.X. Evaluación del contrato real {[}NOMBRE DEL
CONTRATO{]}}\label{x.-evaluaciuxf3n-del-contrato-real-nombre-del-contrato}

\paragraph{Descripción del
contrato}\label{descripciuxf3n-del-contrato-2}

Además de los contratos vulnerables de referencia, se incluyó el
contrato \textbf{{[}NOMBRE{]}}, obtenido de \textbf{{[}REPOSITORIO O
FUENTE{]}}, con el objetivo de analizar el comportamiento de EVMAudit en
un escenario más próximo a un entorno real.

!TODO: describir brevemente la finalidad del contrato.

\begin{center}\rule{0.5\linewidth}{0.5pt}\end{center}

\paragraph{Resultados obtenidos}\label{resultados-obtenidos}

La Tabla X resume las vulnerabilidades detectadas durante el análisis.

{\def\LTcaptype{none} % do not increment counter
\begin{longtable}[]{@{}ll@{}}
\toprule\noalign{}
Métrica & Valor \\
\midrule\noalign{}
\endhead
\bottomrule\noalign{}
\endlastfoot
Vulnerabilidades detectadas por Slither & TODO \\
Vulnerabilidades detectadas por Mythril & TODO \\
Vulnerabilidades finales & TODO \\
Vulnerabilidades confirmadas & TODO \\
\end{longtable}
}

\begin{center}\rule{0.5\linewidth}{0.5pt}\end{center}

\paragraph{Observaciones sobre la
correlación}\label{observaciones-sobre-la-correlaciuxf3n}

En este caso se observó que:

\begin{itemize}
\tightlist
\item
  TODO.
\end{itemize}

Asimismo, se identificaron diferencias entre ambas herramientas en
relación con:

\begin{itemize}
\tightlist
\item
  TODO.
\end{itemize}

\begin{center}\rule{0.5\linewidth}{0.5pt}\end{center}

\paragraph{Resultados del fuzzing}\label{resultados-del-fuzzing-1}

La fase de fuzzing permitió:

\begin{itemize}
\tightlist
\item
  TODO.
\end{itemize}

No obstante, algunas vulnerabilidades no pudieron validarse
automáticamente debido a:

\begin{itemize}
\tightlist
\item
  TODO.
\end{itemize}

\begin{center}\rule{0.5\linewidth}{0.5pt}\end{center}

\paragraph{Discusión}\label{discusiuxf3n-2}

El análisis realizado demuestra que EVMAudit puede aplicarse también
sobre contratos reales y no únicamente sobre ejemplos académicos
diseñados específicamente para contener vulnerabilidades.

!TODO: indicar limitaciones observadas durante el análisis.
